\documentclass[11pt,a4paper]{article}

% =========================================================================
% PAKET-PAKET UTAMA
% =========================================================================
\usepackage[utf8]{inputenc}
\usepackage[T1]{fontenc}
\usepackage{lmodern}
\usepackage[indonesian]{babel}
\usepackage{amsmath,amssymb,amsfonts,amsthm}
\usepackage{bm}              % Notasi tebal untuk vektor & matriks
\usepackage{geometry}
\geometry{top=2.5cm, bottom=2.5cm, left=2.5cm, right=2.5cm}
\usepackage{graphicx}
\usepackage{booktabs}        % Tabel profesional
\usepackage{tabularx}
\usepackage{microtype}
\usepackage{xcolor}
\usepackage{mdframed}        % Kotak highlight untuk definisi
\usepackage{caption}
\usepackage{subcaption}
\usepackage{float}
\usepackage{tikz}
\usetikzlibrary{shapes.geometric, arrows.meta, positioning, calc}
\usepackage{hyperref}
\hypersetup{
    colorlinks=true,
    linkcolor=blue!70!black,
    citecolor=blue!70!black,
    urlcolor=blue!70!black,
    pdfencoding=auto,
    psdextra
}
\pdfstringdefDisableCommands{%
  \def\boldsymbol#1{#1}%
  \def\bm#1{#1}%
  \def\mathbf#1{#1}%
}

% Makro sitasi fleksibel
\newcommand{\citep}[2][]{(Rasmussen \& Williams, 2006\if\relax\detokenize{#1}\relax\else, #1\fi)}
\newcommand{\citet}[2][]{Rasmussen \& Williams (2006\if\relax\detokenize{#1}\relax\else, #1\fi)}

% Lingkungan Teorema & Definisi
\theoremstyle{plain}
\newtheorem{teorema}{Teorema}[section]
\newtheorem{claim}{Claim}[section]

\theoremstyle{definition}
\newtheorem{definisi}{Definisi}[section]

\theoremstyle{remark}
\newtheorem{catatan}{Catatan}[section]

% Format nama pembuktian
\addto\captionsindonesian{\renewcommand{\proofname}{Proof}}
\renewcommand{\proofname}{Proof}

% Path direktori gambar
\graphicspath{{figures/}}

% =========================================================================
% INFORMASI DOKUMEN
% =========================================================================
\title{
    \vspace{-1.2cm}
    \textbf{\Large BIMBINGAN TUGAS AKHIR \#6} \\[0.2em]
}
\author{}
\date{Rabu, 21 Oktober 2026}

\begin{document}

\maketitle
\vspace{-0.8cm}

% =========================================================================
% KONTEN DOKUMEN BIMBINGAN 06
% =========================================================================

% =========================================================================
% 1. PEMBUKTIAN DISTRIBUSI NORMAL y: DUA KASUS BOBOT w
% =========================================================================
\section{Pembuktian Distribusi Normal \texorpdfstring{$y = f(x) + \varepsilon$}{y = f(x) + epsilon}}
\label{sec:pembuktian_mgf}

\noindent Diberikan model observasi linier:
\begin{equation}
y = f(x) + \varepsilon = wx + \varepsilon \label{eq:model_dasar}
\end{equation}
dengan galat $\varepsilon \sim \mathcal{N}(0, \sigma_n^2)$. 

\vspace{0.2cm}
\noindent \textbf{MGF Distribusi Normal}:
\begin{itemize}
    \item \textbf{Univariat}: Untuk peubah acak $Z \sim \mathcal{N}(\mu, \sigma^2)$, MGF didefinisikan sebagai:
    \begin{equation}
    M_Z(t) = \mathbb{E}\left[e^{tZ}\right] = \exp\left( \mu t + \frac{1}{2}\sigma^2 t^2 \right) \label{eq:mgf_normal_univariat}
    \end{equation}
    \item \textbf{Multivariat}: Untuk vektor acak $\mathbf{Z} \sim \mathcal{N}(\boldsymbol{\mu}, \Sigma)$ dengan $\mathbf{Z} \in \mathbb{R}^k$, MGF didefinisikan sebagai:
    \begin{equation}
    M_{\mathbf{Z}}(\mathbf{t}) = \mathbb{E}\left[\exp(\mathbf{t}^T \mathbf{Z})\right] = \exp\left( \mathbf{t}^T \boldsymbol{\mu} + \frac{1}{2}\mathbf{t}^T \Sigma \mathbf{t} \right) \label{eq:mgf_normal_multivariat}
    \end{equation}
\end{itemize}
Untuk galat $\varepsilon \sim \mathcal{N}(0, \sigma_n^2)$, MGF-nya adalah $M_\varepsilon(t) = \exp\left( \frac{1}{2}\sigma_n^2 t^2 \right)$.

Pembuktian distribusi observasi ditinjau dalam dua kasus nilai bobot:
\begin{enumerate}
    \item \textbf{Kasus 1}: Bobot bernilai konstan ($w \in \mathbb{R}$ atau $\mathbf{w} \in \mathbb{R}^d$).
    \item \textbf{Kasus 2}: Bobot berdistribusi normal ($w \sim \mathcal{N}(0, \sigma_w^2)$ atau $\mathbf{w} \sim \mathcal{N}(\mathbf{0}, \Sigma_p)$).
\end{enumerate}

% -------------------------------------------------------------------------
\subsection{Kasus 1: Bobot Bernilai Konstan}
\label{subsec:kasus_1}

\begin{claim}[Bentuk Skalar Univariat]
Jika $w \in \mathbb{R}$ berupa skalar konstan dan galat $\varepsilon \sim \mathcal{N}(0, \sigma_n^2)$, maka variabel observasi $y = wx + \varepsilon$ berdistribusi normal:
\begin{equation}
y \sim \mathcal{N}\left( wx, \; \sigma_n^2 \right)
\end{equation}
\end{claim}

\begin{proof}
MGF dari $y$:
\begin{align}
M_y(t) &= \mathbb{E}\left[ e^{t(wx + \varepsilon)} \right] \nonumber \\
&= \mathbb{E}\left[ e^{twx} e^{t\varepsilon} \right] \nonumber \\
&= e^{twx} \, \mathbb{E}\left[ e^{t\varepsilon} \right] \nonumber \\
&= e^{twx} M_\varepsilon(t) \nonumber \\
&= e^{twx} \exp\left( \frac{1}{2}\sigma_n^2 t^2 \right) \nonumber \\
&= \exp\left( (wx)t + \frac{1}{2}\sigma_n^2 t^2 \right) \label{eq:mgf_y_kasus1}
\end{align}
Berdasarkan bentuk MGF normal univariat pada Persamaan \eqref{eq:mgf_normal_univariat} dengan $\mu = wx$ dan $\sigma^2 = \sigma_n^2$, terbukti bahwa:
\begin{equation}
y \sim \mathcal{N}\left( wx, \; \sigma_n^2 \right) \qedhere
\end{equation}
\end{proof}

\begin{claim}[Bentuk Vektor Multivariat: Persamaan (2.3) GPML]
\label{claim:likelihood_gpml}
Untuk sekumpulan input $X = [\mathbf{x}_1, \dots, \mathbf{x}_n] \in \mathbb{R}^{d \times n}$, bobot $\mathbf{w} \in \mathbb{R}^d$ bernilai konstan, dan galat $\boldsymbol{\varepsilon} \sim \mathcal{N}(\mathbf{0}, \sigma_n^2 I_n)$ \citep[hlm.~9]{rasmussen2006gaussian}, distribusi bersyarat observasi $\mathbf{y}$ berdistribusi normal multivariat:
\begin{equation}
\mathbf{y} \mid X, \mathbf{w} \sim \mathcal{N}\left( X^T \mathbf{w}, \; \sigma_n^2 I_n \right)
\end{equation}
dengan fungsi kepadatan probabilitas (likelihood) yang terfaktorisasi:
\begin{equation}
p(\mathbf{y} \mid X, \mathbf{w}) = \prod_{i=1}^n p(y_i \mid \mathbf{x}_i, \mathbf{w}) = \frac{1}{(2\pi \sigma_n^2)^{n/2}} \exp\left( -\frac{1}{2\sigma_n^2} |\mathbf{y} - X^T \mathbf{w}|^2 \right) \label{eq:gpml_eq_2_3}
\end{equation}
\end{claim}

\begin{proof}
MGF bersyarat dari $\mathbf{y} \mid X, \mathbf{w}$ didefinisikan sebagai:
\begin{align}
M_{\mathbf{y} \mid X, \mathbf{w}}(\mathbf{t}) &= \mathbb{E}\left[ \exp\left( \mathbf{t}^T (X^T \mathbf{w} + \boldsymbol{\varepsilon}) \right) \mid X, \mathbf{w} \right] \nonumber \\
&= \exp\left( \mathbf{t}^T (X^T \mathbf{w}) \right) \cdot \mathbb{E}\left[ \exp\left( \mathbf{t}^T \boldsymbol{\varepsilon} \right) \right] \nonumber \\
&= \exp\left( \mathbf{t}^T (X^T \mathbf{w}) \right) \cdot M_{\boldsymbol{\varepsilon}}(\mathbf{t}) \nonumber \\
&= \exp\left( \mathbf{t}^T (X^T \mathbf{w}) + \frac{1}{2}\mathbf{t}^T \left( \sigma_n^2 I_n \right) \mathbf{t} \right) \label{eq:mgf_likelihood_multivariat}
\end{align}
Berdasarkan MGF normal multivariat pada Persamaan \eqref{eq:mgf_normal_multivariat} dengan $\boldsymbol{\mu} = X^T \mathbf{w}$ dan $\Sigma = \sigma_n^2 I_n$, terbukti bahwa $\mathbf{y} \mid X, \mathbf{w} \sim \mathcal{N}\left( X^T \mathbf{w}, \; \sigma_n^2 I_n \right)$.

Karena $\Sigma = \sigma_n^2 I_n$ diagonal, suku eksponen pada MGF \eqref{eq:mgf_likelihood_multivariat} dapat diuraikan:
\begin{equation}
M_{\mathbf{y} \mid X, \mathbf{w}}(\mathbf{t}) = \prod_{i=1}^n \exp\left( t_i (\mathbf{x}_i^T \mathbf{w}) + \frac{1}{2}\sigma_n^2 t_i^2 \right) = \prod_{i=1}^n M_{y_i \mid \mathbf{x}_i, \mathbf{w}}(t_i)
\end{equation}
Karena MGF gabungan bersyarat sama dengan hasil kali MGF masing-masing variabel, maka setiap $y_1, \dots, y_n$ saling bebas dengan $y_i \mid \mathbf{x}_i, \mathbf{w} \sim \mathcal{N}(\mathbf{x}_i^T \mathbf{w}, \sigma_n^2)$. Oleh karena itu, fungsi kepadatan bersama terfaktorisasi menjadi:
\begin{equation}
p(\mathbf{y} \mid X, \mathbf{w}) = \prod_{i=1}^n p(y_i \mid \mathbf{x}_i, \mathbf{w}) = \frac{1}{(2\pi \sigma_n^2)^{n/2}} \exp\left( -\frac{1}{2\sigma_n^2} |\mathbf{y} - X^T \mathbf{w}|^2 \right) \qedhere
\end{equation}
\end{proof}

% -------------------------------------------------------------------------
\subsection{Kasus 2: Bobot \texorpdfstring{$w$}{w} Berdistribusi Normal (\texorpdfstring{$w \sim \mathcal{N}(0, \sigma_w^2)$}{w Gaussian})}
\label{subsec:kasus_2}

Diberikan $w \sim \mathcal{N}(0, \sigma_w^2)$ dan diasumsikan $w$ dan $\varepsilon$ saling bebas.

\begin{claim}
Fungsi $f(x) = wx$ berdistribusi normal:
\begin{equation}
f(x) \sim \mathcal{N}\left( 0, \; x^2 \sigma_w^2 \right)
\end{equation}
\end{claim}

\begin{proof}
MGF dari $w$ adalah $M_w(t) = \exp\left( \frac{1}{2}\sigma_w^2 t^2 \right)$. MGF dari $f(x) = wx$:
\begin{align}
M_{f(x)}(t) &= \mathbb{E}_w\left[ e^{t(wx)} \right] = \mathbb{E}_w\left[ e^{(tx)w} \right] = M_w(tx) \nonumber \\
&= \exp\left( \frac{1}{2}\sigma_w^2 (tx)^2 \right) \nonumber \\
&= \exp\left( \frac{1}{2} t^2 \left( x^2 \sigma_w^2 \right) \right) \label{eq:mgf_fx_kasus2}
\end{align}
Sesuai bentuk MGF normal univariat pada Persamaan \eqref{eq:mgf_normal_univariat} dengan $\mu = 0$ dan variansi $\sigma^2 = x^2 \sigma_w^2$, terbukti bahwa:
\begin{equation}
f(x) \sim \mathcal{N}\left( 0, \; x^2 \sigma_w^2 \right) \qedhere
\end{equation}
\end{proof}

\begin{claim}
Variabel observasi $y = f(x) + \varepsilon$ berdistribusi normal:
\begin{equation}
y \sim \mathcal{N}\left( 0, \; x^2 \sigma_w^2 + \sigma_n^2 \right)
\end{equation}
\end{claim}

\begin{proof}
Karena $w$ dan $\varepsilon$ saling bebas, maka $f(x)$ dan $\varepsilon$ saling bebas. MGF dari $y = f(x) + \varepsilon$:
\begin{align}
M_y(t) &= \mathbb{E}_{w, \varepsilon}\left[ e^{t(f(x) + \varepsilon)} \right] \nonumber \\
&= \mathbb{E}\left[ e^{t f(x)} \right] \cdot \mathbb{E}\left[ e^{t\varepsilon} \right] \nonumber \\
&= M_{f(x)}(t) \cdot M_\varepsilon(t) \nonumber \\
&= \exp\left( \frac{1}{2} t^2 \left( x^2 \sigma_w^2 \right) \right) \cdot \exp\left( \frac{1}{2} \sigma_n^2 t^2 \right) \nonumber \\
&= \exp\left( \frac{1}{2} t^2 \left( x^2 \sigma_w^2 + \sigma_n^2 \right) \right) \label{eq:mgf_y_kasus2}
\end{align}
Sesuai bentuk MGF normal univariat pada Persamaan \eqref{eq:mgf_normal_univariat} dengan $\mu = 0$ dan variansi $\sigma^2 = x^2 \sigma_w^2 + \sigma_n^2$, terbukti bahwa:
\begin{equation}
y \sim \mathcal{N}\left( 0, \; x^2 \sigma_w^2 + \sigma_n^2 \right) \qedhere
\end{equation}
\end{proof}

\begin{claim}[Bentuk Vektor Multivariat]
Untuk sekumpulan input $X = [\mathbf{x}_1, \dots, \mathbf{x}_n]^T \in \mathbb{R}^{n \times d}$, bobot $\mathbf{w} \sim \mathcal{N}(\mathbf{0}, \Sigma_p)$, dan galat $\boldsymbol{\varepsilon} \sim \mathcal{N}(\mathbf{0}, \sigma_n^2 I_n)$ dengan $\mathbf{w}$ dan $\boldsymbol{\varepsilon}$ saling bebas:
\begin{equation}
\mathbf{f} \sim \mathcal{N}\left(\mathbf{0}, \; X\Sigma_p X^T\right) \quad \text{dan} \quad \mathbf{y} \sim \mathcal{N}\left(\mathbf{0}, \; X\Sigma_p X^T + \sigma_n^2 I_n\right)
\end{equation}
\end{claim}

\begin{proof}
Berdasarkan MGF normal multivariat pada Persamaan \eqref{eq:mgf_normal_multivariat}:
\begin{equation}
M_{\mathbf{f}}(\mathbf{t}) = \mathbb{E}\left[ \exp(\mathbf{t}^T X \mathbf{w}) \right] = \exp\left( \frac{1}{2}\mathbf{t}^T \left( X\Sigma_p X^T \right) \mathbf{t} \right) \implies \mathbf{f} \sim \mathcal{N}\left(\mathbf{0}, \; X\Sigma_p X^T\right)
\end{equation}
Karena $\mathbf{w}$ dan $\boldsymbol{\varepsilon}$ saling bebas:
\begin{align}
M_{\mathbf{y}}(\mathbf{t}) &= M_{\mathbf{f}}(\mathbf{t}) \cdot M_{\boldsymbol{\varepsilon}}(\mathbf{t}) \nonumber \\
&= \exp\left( \frac{1}{2}\mathbf{t}^T \left( X\Sigma_p X^T + \sigma_n^2 I_n \right) \mathbf{t} \right) \implies \mathbf{y} \sim \mathcal{N}\left(\mathbf{0}, \; X\Sigma_p X^T + \sigma_n^2 I_n\right) \qedhere
\end{align}
\end{proof}

% =========================================================================
% 2. PERBEDAAN DNN DAN DGP
% =========================================================================
\newpage
\section{Perbedaan DNN dan DGP}
\label{sec:perbedaan_dnn_dgp}

\noindent \textbf{Deep Neural Network (DNN)}:
\begin{align}
\mathbf{h}^{(0)} &= \mathbf{x} \in \mathbb{R}^{d_0} \nonumber \\
\mathbf{h}^{(\ell)} &= \sigma\left( W^{(\ell)} \mathbf{h}^{(\ell-1)} + \mathbf{b}^{(\ell)} \right), \quad \ell = 1, \dots, L \label{eq:dnn_layer} \\
\hat{\mathbf{y}} &= \mathbf{h}^{(L)} = (f_L \circ \dots \circ f_1)(\mathbf{x}) \label{eq:dnn_composition}
\end{align}
dengan parameter bobot deterministik $W^{(\ell)} \in \mathbb{R}^{d_\ell \times d_{\ell-1}}$, bias $\mathbf{b}^{(\ell)} \in \mathbb{R}^{d_\ell}$, dan fungsi aktivasi non-linier $\sigma(\cdot)$.

\vspace{0.2cm}
\noindent \textbf{Gaussian Process (GP) Standar}:
\begin{equation}
f(\mathbf{x}) \sim \mathcal{GP}\left( 0, \; k(\mathbf{x}, \mathbf{x}') \right), \quad y = f(\mathbf{x}) + \varepsilon, \quad \varepsilon \sim \mathcal{N}\left(0, \; \sigma_n^2\right) \label{eq:gp_standard}
\end{equation}
dengan fungsi kovariansi tunggal $k(\mathbf{x}, \mathbf{x}')$.

\vspace{0.2cm}
\noindent \textbf{Deep Gaussian Process (DGP)}:
\begin{align}
\mathbf{h}^{(0)} &= \mathbf{x} \nonumber \\
\mathbf{h}^{(\ell)} &= \mathbf{f}^{(\ell)}(\mathbf{h}^{(\ell-1)}) + \boldsymbol{\varepsilon}^{(\ell)}, \quad \ell = 1, \dots, L \label{eq:dgp_layer}
\end{align}
di mana pemetaan berbobot $(W^{(\ell)}, \mathbf{b}^{(\ell)})$ pada \eqref{eq:dnn_layer} digantikan oleh fungsi dari distribusi proses Gaussian ($\mathcal{GP}$):
\begin{equation}
\mathbf{f}^{(\ell)} \sim \mathcal{GP}\left( \mathbf{0}, \; K^{(\ell)} \right), \quad K^{(\ell)}_{ij} = k^{(\ell)}\left(\mathbf{h}^{(\ell-1)}_i, \mathbf{h}^{(\ell-1)}_j\right), \quad \boldsymbol{\varepsilon}^{(\ell)} \sim \mathcal{N}\left(\mathbf{0}, \; \sigma_\ell^2 I\right) \label{eq:dgp_prior}
\end{equation}
Komposisi \textit{output} akhir:
\begin{equation}
\mathbf{y} = \left(\mathbf{f}^{(L)} \circ \dots \circ \mathbf{f}^{(1)}\right)(\mathbf{x}) + \boldsymbol{\varepsilon} = \mathbf{f}^{(L)}\left( \dots \mathbf{f}^{(1)}(\mathbf{x}) \dots \right) + \boldsymbol{\varepsilon} \label{eq:dgp_composition}
\end{equation}

\end{document}
